% Options for packages loaded elsewhere
\PassOptionsToPackage{unicode}{hyperref}
\PassOptionsToPackage{hyphens}{url}
%
\documentclass[
]{article}
\usepackage{amsmath,amssymb}
\usepackage{iftex}
\ifPDFTeX
  \usepackage[T1]{fontenc}
  \usepackage[utf8]{inputenc}
  \usepackage{textcomp} % provide euro and other symbols
\else % if luatex or xetex
  \usepackage{unicode-math} % this also loads fontspec
  \defaultfontfeatures{Scale=MatchLowercase}
  \defaultfontfeatures[\rmfamily]{Ligatures=TeX,Scale=1}
\fi
\usepackage{lmodern}
\ifPDFTeX\else
  % xetex/luatex font selection
\fi
% Use upquote if available, for straight quotes in verbatim environments
\IfFileExists{upquote.sty}{\usepackage{upquote}}{}
\IfFileExists{microtype.sty}{% use microtype if available
  \usepackage[]{microtype}
  \UseMicrotypeSet[protrusion]{basicmath} % disable protrusion for tt fonts
}{}
\makeatletter
\@ifundefined{KOMAClassName}{% if non-KOMA class
  \IfFileExists{parskip.sty}{%
    \usepackage{parskip}
  }{% else
    \setlength{\parindent}{0pt}
    \setlength{\parskip}{6pt plus 2pt minus 1pt}}
}{% if KOMA class
  \KOMAoptions{parskip=half}}
\makeatother
\usepackage{xcolor}
\usepackage[margin=1in]{geometry}
\usepackage{color}
\usepackage{fancyvrb}
\newcommand{\VerbBar}{|}
\newcommand{\VERB}{\Verb[commandchars=\\\{\}]}
\DefineVerbatimEnvironment{Highlighting}{Verbatim}{commandchars=\\\{\}}
% Add ',fontsize=\small' for more characters per line
\usepackage{framed}
\definecolor{shadecolor}{RGB}{248,248,248}
\newenvironment{Shaded}{\begin{snugshade}}{\end{snugshade}}
\newcommand{\AlertTok}[1]{\textcolor[rgb]{0.94,0.16,0.16}{#1}}
\newcommand{\AnnotationTok}[1]{\textcolor[rgb]{0.56,0.35,0.01}{\textbf{\textit{#1}}}}
\newcommand{\AttributeTok}[1]{\textcolor[rgb]{0.13,0.29,0.53}{#1}}
\newcommand{\BaseNTok}[1]{\textcolor[rgb]{0.00,0.00,0.81}{#1}}
\newcommand{\BuiltInTok}[1]{#1}
\newcommand{\CharTok}[1]{\textcolor[rgb]{0.31,0.60,0.02}{#1}}
\newcommand{\CommentTok}[1]{\textcolor[rgb]{0.56,0.35,0.01}{\textit{#1}}}
\newcommand{\CommentVarTok}[1]{\textcolor[rgb]{0.56,0.35,0.01}{\textbf{\textit{#1}}}}
\newcommand{\ConstantTok}[1]{\textcolor[rgb]{0.56,0.35,0.01}{#1}}
\newcommand{\ControlFlowTok}[1]{\textcolor[rgb]{0.13,0.29,0.53}{\textbf{#1}}}
\newcommand{\DataTypeTok}[1]{\textcolor[rgb]{0.13,0.29,0.53}{#1}}
\newcommand{\DecValTok}[1]{\textcolor[rgb]{0.00,0.00,0.81}{#1}}
\newcommand{\DocumentationTok}[1]{\textcolor[rgb]{0.56,0.35,0.01}{\textbf{\textit{#1}}}}
\newcommand{\ErrorTok}[1]{\textcolor[rgb]{0.64,0.00,0.00}{\textbf{#1}}}
\newcommand{\ExtensionTok}[1]{#1}
\newcommand{\FloatTok}[1]{\textcolor[rgb]{0.00,0.00,0.81}{#1}}
\newcommand{\FunctionTok}[1]{\textcolor[rgb]{0.13,0.29,0.53}{\textbf{#1}}}
\newcommand{\ImportTok}[1]{#1}
\newcommand{\InformationTok}[1]{\textcolor[rgb]{0.56,0.35,0.01}{\textbf{\textit{#1}}}}
\newcommand{\KeywordTok}[1]{\textcolor[rgb]{0.13,0.29,0.53}{\textbf{#1}}}
\newcommand{\NormalTok}[1]{#1}
\newcommand{\OperatorTok}[1]{\textcolor[rgb]{0.81,0.36,0.00}{\textbf{#1}}}
\newcommand{\OtherTok}[1]{\textcolor[rgb]{0.56,0.35,0.01}{#1}}
\newcommand{\PreprocessorTok}[1]{\textcolor[rgb]{0.56,0.35,0.01}{\textit{#1}}}
\newcommand{\RegionMarkerTok}[1]{#1}
\newcommand{\SpecialCharTok}[1]{\textcolor[rgb]{0.81,0.36,0.00}{\textbf{#1}}}
\newcommand{\SpecialStringTok}[1]{\textcolor[rgb]{0.31,0.60,0.02}{#1}}
\newcommand{\StringTok}[1]{\textcolor[rgb]{0.31,0.60,0.02}{#1}}
\newcommand{\VariableTok}[1]{\textcolor[rgb]{0.00,0.00,0.00}{#1}}
\newcommand{\VerbatimStringTok}[1]{\textcolor[rgb]{0.31,0.60,0.02}{#1}}
\newcommand{\WarningTok}[1]{\textcolor[rgb]{0.56,0.35,0.01}{\textbf{\textit{#1}}}}
\usepackage{longtable,booktabs,array}
\usepackage{calc} % for calculating minipage widths
% Correct order of tables after \paragraph or \subparagraph
\usepackage{etoolbox}
\makeatletter
\patchcmd\longtable{\par}{\if@noskipsec\mbox{}\fi\par}{}{}
\makeatother
% Allow footnotes in longtable head/foot
\IfFileExists{footnotehyper.sty}{\usepackage{footnotehyper}}{\usepackage{footnote}}
\makesavenoteenv{longtable}
\usepackage{graphicx}
\makeatletter
\def\maxwidth{\ifdim\Gin@nat@width>\linewidth\linewidth\else\Gin@nat@width\fi}
\def\maxheight{\ifdim\Gin@nat@height>\textheight\textheight\else\Gin@nat@height\fi}
\makeatother
% Scale images if necessary, so that they will not overflow the page
% margins by default, and it is still possible to overwrite the defaults
% using explicit options in \includegraphics[width, height, ...]{}
\setkeys{Gin}{width=\maxwidth,height=\maxheight,keepaspectratio}
% Set default figure placement to htbp
\makeatletter
\def\fps@figure{htbp}
\makeatother
\setlength{\emergencystretch}{3em} % prevent overfull lines
\providecommand{\tightlist}{%
  \setlength{\itemsep}{0pt}\setlength{\parskip}{0pt}}
\setcounter{secnumdepth}{-\maxdimen} % remove section numbering
\usepackage{booktabs}
\usepackage{longtable}
\usepackage{array}
\usepackage{multirow}
\usepackage{wrapfig}
\usepackage{float}
\usepackage{colortbl}
\usepackage{pdflscape}
\usepackage{tabu}
\usepackage{threeparttable}
\usepackage{threeparttablex}
\usepackage[normalem]{ulem}
\usepackage{makecell}
\usepackage{xcolor}
\ifLuaTeX
  \usepackage{selnolig}  % disable illegal ligatures
\fi
\usepackage{bookmark}
\IfFileExists{xurl.sty}{\usepackage{xurl}}{} % add URL line breaks if available
\urlstyle{same}
\hypersetup{
  pdftitle={PAC\_2},
  hidelinks,
  pdfcreator={LaTeX via pandoc}}

\title{PAC\_2}
\author{}
\date{\vspace{-2.5em}2025-04-20}

\begin{document}
\maketitle

\subsection{Breu anàlisi visual}\label{breu-anuxe0lisi-visual}

A partir de les preguntes de recerca formulades, representeu visualment
les variables a estudiar de manera que orientin les respostes a aquestes
preguntes. És a dir, per cada pregunta de recerca, representeu un gràfic
que ajudi a interpretar visualment la distribució de les variables per
aquella pregunta. Escolliu el gràfic que sigui més apropiat en funció
del vostre criteri. Estructureu aquesta secció en diferents subseccions,
una per cada pregunta de recerca.

{Llegir el fitxer de dades sleephealth.csv.}

\begin{Shaded}
\begin{Highlighting}[]
\FunctionTok{library}\NormalTok{(readxl)}
\NormalTok{data }\OtherTok{\textless{}{-}} \FunctionTok{read\_excel}\NormalTok{(}\StringTok{"\textasciitilde{}/Documentos/GitHub/UOC{-}treballs/analisis\_estadistic/pac2/sleephealth\_processed.xlsx"}\NormalTok{)}
\FunctionTok{kable}\NormalTok{(}\FunctionTok{head}\NormalTok{(data))}
\end{Highlighting}
\end{Shaded}

\begin{longtable}[]{@{}
  >{\raggedleft\arraybackslash}p{(\columnwidth - 28\tabcolsep) * \real{0.0291}}
  >{\raggedright\arraybackslash}p{(\columnwidth - 28\tabcolsep) * \real{0.0680}}
  >{\raggedleft\arraybackslash}p{(\columnwidth - 28\tabcolsep) * \real{0.0388}}
  >{\raggedright\arraybackslash}p{(\columnwidth - 28\tabcolsep) * \real{0.1748}}
  >{\raggedleft\arraybackslash}p{(\columnwidth - 28\tabcolsep) * \real{0.0388}}
  >{\raggedleft\arraybackslash}p{(\columnwidth - 28\tabcolsep) * \real{0.0291}}
  >{\raggedleft\arraybackslash}p{(\columnwidth - 28\tabcolsep) * \real{0.0388}}
  >{\raggedleft\arraybackslash}p{(\columnwidth - 28\tabcolsep) * \real{0.0680}}
  >{\raggedright\arraybackslash}p{(\columnwidth - 28\tabcolsep) * \real{0.1068}}
  >{\raggedright\arraybackslash}p{(\columnwidth - 28\tabcolsep) * \real{0.0680}}
  >{\raggedleft\arraybackslash}p{(\columnwidth - 28\tabcolsep) * \real{0.0291}}
  >{\raggedleft\arraybackslash}p{(\columnwidth - 28\tabcolsep) * \real{0.0583}}
  >{\raggedright\arraybackslash}p{(\columnwidth - 28\tabcolsep) * \real{0.1165}}
  >{\raggedleft\arraybackslash}p{(\columnwidth - 28\tabcolsep) * \real{0.0680}}
  >{\raggedleft\arraybackslash}p{(\columnwidth - 28\tabcolsep) * \real{0.0680}}@{}}
\toprule\noalign{}
\begin{minipage}[b]{\linewidth}\raggedleft
ID
\end{minipage} & \begin{minipage}[b]{\linewidth}\raggedright
Gender
\end{minipage} & \begin{minipage}[b]{\linewidth}\raggedleft
Age
\end{minipage} & \begin{minipage}[b]{\linewidth}\raggedright
Occupation
\end{minipage} & \begin{minipage}[b]{\linewidth}\raggedleft
SH
\end{minipage} & \begin{minipage}[b]{\linewidth}\raggedleft
SQ
\end{minipage} & \begin{minipage}[b]{\linewidth}\raggedleft
PHY
\end{minipage} & \begin{minipage}[b]{\linewidth}\raggedleft
Stress
\end{minipage} & \begin{minipage}[b]{\linewidth}\raggedright
BMI
\end{minipage} & \begin{minipage}[b]{\linewidth}\raggedright
BP
\end{minipage} & \begin{minipage}[b]{\linewidth}\raggedleft
HR
\end{minipage} & \begin{minipage}[b]{\linewidth}\raggedleft
Steps
\end{minipage} & \begin{minipage}[b]{\linewidth}\raggedright
SD
\end{minipage} & \begin{minipage}[b]{\linewidth}\raggedleft
BPsyst
\end{minipage} & \begin{minipage}[b]{\linewidth}\raggedleft
BPdias
\end{minipage} \\
\midrule\noalign{}
\endhead
\bottomrule\noalign{}
\endlastfoot
1 & Male & 27 & Software Engineer & 6.1 & 6 & 42 & 6 & Overweight &
126/83 & 77 & 4200 & None & 126 & 83 \\
2 & Male & 28 & Doctor & 6.2 & 6 & 60 & 8 & Normal & 125/80 & 75 & 10000
& None & 125 & 80 \\
3 & Male & 28 & Doctor & 6.2 & 6 & 60 & 8 & Normal & 125/80 & 75 & 10000
& None & 125 & 80 \\
4 & Male & 28 & Sales Person & 5.9 & 4 & 30 & 8 & Obese & 140/90 & 85 &
3000 & Sleep Apnea & 140 & 90 \\
5 & Male & 28 & Sales Person & 5.9 & 4 & 30 & 8 & Obese & 140/90 & 85 &
3000 & Sleep Apnea & 140 & 90 \\
6 & Male & 28 & Software Engineer & 5.9 & 4 & 30 & 8 & Obese & 140/90 &
85 & 3000 & Insomnia & 140 & 90 \\
\end{longtable}

\subsubsection{PR1: La mitjana d'hores de son de les persones a l'estudi
és significativament inferior a 7.5 hores
diàries?}\label{pr1-la-mitjana-dhores-de-son-de-les-persones-a-lestudi-uxe9s-significativament-inferior-a-7.5-hores-diuxe0ries}

Farem un histograma de la distribucio ja que ens pot donar una idea de
com esta distribuida i on esta la mitjana

\begin{Shaded}
\begin{Highlighting}[]
\FunctionTok{library}\NormalTok{(ggplot2)}
\FunctionTok{ggplot}\NormalTok{(data, }\FunctionTok{aes}\NormalTok{(}\AttributeTok{x=}\NormalTok{SH)) }\SpecialCharTok{+} \FunctionTok{geom\_histogram}\NormalTok{(}\AttributeTok{bins =} \DecValTok{20}\NormalTok{,}\AttributeTok{color=}\StringTok{"black"}\NormalTok{,}\AttributeTok{fill=}\StringTok{"white"}\NormalTok{) }\SpecialCharTok{+}
  \FunctionTok{labs}\NormalTok{(}\AttributeTok{title =} \StringTok{"Histogram of hours sleep"}\NormalTok{,}
       \AttributeTok{x =} \StringTok{"Hours Sleep"}\NormalTok{,}
       \AttributeTok{y =} \StringTok{"Count"}\NormalTok{) }\SpecialCharTok{+}
  \FunctionTok{theme\_minimal}\NormalTok{()}
\end{Highlighting}
\end{Shaded}

\includegraphics{PAC_2_Albert_lozano_files/figure-latex/unnamed-chunk-1-1.pdf}
Mirant el histograma la distribucio no sembla normal i encara que el
valor mes comú esta prop del 7.5 sent mes o menys 7.3 i tenim valors al
dos costats d'aquest i sembla que tenim una mica més en els valors més
petits per tant sembla que la mitjana estara per sota, pero no ho podem
asegurar fins fer el test estadistic \#\#\# PR2: Existeixen diferències
en la qualitat del son entre gèneres? Farem dos boxplot de la qualitat
del son on separarem els dos generes per veure si hi ha diferencies.

\begin{Shaded}
\begin{Highlighting}[]
\FunctionTok{library}\NormalTok{(ggplot2)}
\FunctionTok{ggplot}\NormalTok{(data, }\FunctionTok{aes}\NormalTok{(}\AttributeTok{x=}\FunctionTok{factor}\NormalTok{(SQ), }\AttributeTok{fill=}\NormalTok{Gender)) }\SpecialCharTok{+}
  \FunctionTok{geom\_bar}\NormalTok{(}\AttributeTok{stat=}\StringTok{"count"}\NormalTok{, }\AttributeTok{width=}\FloatTok{0.7}\NormalTok{, }\AttributeTok{position=}\FunctionTok{position\_dodge}\NormalTok{()) }\SpecialCharTok{+}
  \FunctionTok{labs}\NormalTok{(}\AttributeTok{title =} \StringTok{"Comparison of sleep quality by gender"}\NormalTok{,}
       \AttributeTok{x =} \StringTok{"Sleep Quality"}\NormalTok{,}
       \AttributeTok{y =} \StringTok{"Count"}\NormalTok{) }\SpecialCharTok{+}
  \FunctionTok{theme\_minimal}\NormalTok{()}
\end{Highlighting}
\end{Shaded}

\includegraphics{PAC_2_Albert_lozano_files/figure-latex/unnamed-chunk-2-1.pdf}
Sembla que les dones tenen una qualitat del son superior ja que tenen
més valors més alts \#\#\# PR3: La proporció de persones amb trastorn
del son és més gran en les persones amb nivells alts d'estrès en
comparació amb les que presenten nivells baixos d'estrès? Farem tres
boxplot de l'estres on separarem els dos generes per veure si hi ha
diferencies.

\begin{Shaded}
\begin{Highlighting}[]
\FunctionTok{library}\NormalTok{(ggplot2)}
\FunctionTok{ggplot}\NormalTok{(data, }\FunctionTok{aes}\NormalTok{(}\AttributeTok{x=}\FunctionTok{factor}\NormalTok{(Stress), }\AttributeTok{fill=}\NormalTok{SD)) }\SpecialCharTok{+}
  \FunctionTok{geom\_bar}\NormalTok{(}\AttributeTok{stat=}\StringTok{"count"}\NormalTok{, }\AttributeTok{width=}\FloatTok{0.7}\NormalTok{, }\AttributeTok{position=}\FunctionTok{position\_dodge}\NormalTok{()) }\SpecialCharTok{+}
  \FunctionTok{labs}\NormalTok{(}\AttributeTok{title =} \StringTok{"Comparison of stress by sleep disorders"}\NormalTok{,}
       \AttributeTok{x =} \StringTok{"Stress"}\NormalTok{,}
       \AttributeTok{y =} \StringTok{"Count"}\NormalTok{) }\SpecialCharTok{+}
  \FunctionTok{theme\_minimal}\NormalTok{()}
\end{Highlighting}
\end{Shaded}

\includegraphics{PAC_2_Albert_lozano_files/figure-latex/unnamed-chunk-3-1.pdf}
Sembla que les persones sense trastorns del son tenen uns nivells
d'estres corresponents a una distribució normal mentre que els que tenen
algun trastorn del son tenen mes pics als extrems però els pics els
tenen tant en baix com alt estres per tant es difícil dir si en general
tenen mes estres

\subsubsection{PR4: Hi ha una relació significativa entre la categoria
de BMI (body mass index) i la presència de trastorns del
son?}\label{pr4-hi-ha-una-relaciuxf3-significativa-entre-la-categoria-de-bmi-body-mass-index-i-la-presuxe8ncia-de-trastorns-del-son}

Farem tres boxplot de el nombre de persones amb desordres del son on
separarem els dos generes per veure si hi ha diferencies.

\begin{Shaded}
\begin{Highlighting}[]
\FunctionTok{library}\NormalTok{(ggplot2)}
\FunctionTok{ggplot}\NormalTok{(data, }\FunctionTok{aes}\NormalTok{(}\AttributeTok{x=}\FunctionTok{factor}\NormalTok{(SD), }\AttributeTok{fill=}\NormalTok{BMI)) }\SpecialCharTok{+}
  \FunctionTok{geom\_bar}\NormalTok{(}\AttributeTok{stat=}\StringTok{"count"}\NormalTok{, }\AttributeTok{width=}\FloatTok{0.7}\NormalTok{, }\AttributeTok{position=}\FunctionTok{position\_dodge}\NormalTok{()) }\SpecialCharTok{+}
  \FunctionTok{labs}\NormalTok{(}\AttributeTok{title =} \StringTok{"Comparison of sleep disorders by BMI"}\NormalTok{,}
       \AttributeTok{x =} \StringTok{"Sleep Disorders"}\NormalTok{,}
       \AttributeTok{y =} \StringTok{"Count"}\NormalTok{) }\SpecialCharTok{+}
  \FunctionTok{theme\_minimal}\NormalTok{()}
\end{Highlighting}
\end{Shaded}

\includegraphics{PAC_2_Albert_lozano_files/figure-latex/unnamed-chunk-4-1.pdf}
Sembla que si ja que les persones amb sobrepes estan més representades
als desordres del son comparat amb els que no \#\# Mitjana d'hores de
son (PR1) En aquesta secció volem respondre la pregunta de recerca: PR1:
La mitjana d'hores de son de la població és significativament inferior a
7.5 hores diàries? Per a donar-hi resposta, seguiu els passos que
s'indiquen a continuació \#\#\# Hipòtesis

\end{document}
